\documentclass[11pt]{article}

\usepackage[T1]{fontenc}
\usepackage{amsmath,amssymb,amsthm}
\usepackage[margin=1in]{geometry}
\usepackage[colorlinks=true,linkcolor=blue,urlcolor=blue]{hyperref}

\newtheorem{theorem}{Theorem}
\newtheorem{proposition}[theorem]{Proposition}
\newtheorem{corollary}[theorem]{Corollary}
\newtheorem{conjecture}[theorem]{Conjecture}
\theoremstyle{remark}
\newtheorem{remark}[theorem]{Remark}

\newcommand{\R}{\mathbb R}
\newcommand{\Pp}{\mathbb P}
\newcommand{\norm}[1]{\left\|#1\right\|}
\newcommand{\ip}[2]{\left\langle #1,#2\right\rangle}
\newcommand{\weakto}{\rightharpoonup}
\setlength{\emergencystretch}{2em}

\title{The q4 Nonlinear Entrance Forces\\
a Critical Boundary Heat Impulse}
\author{John Lawson and Codex}
\date{25 July 2026}

\begin{document}
\maketitle

\begin{abstract}
Inside the conditional energy-efficient exact q4 branch, the preceding
round converted the full terminal kinetic defect into a forward
perturbed Navier--Stokes entrance \(v\).  It has weak vector trace zero
but positive limiting energy \(d\).  This round passes its positive-time
mild equation all the way to that boundary.  The fixed-time heat
evolution of the weak-zero datum vanishes, giving
\[
 v(\tau)
 =
 -\int_0^\tau e^{\nu(\tau-s)\Delta}
 \mathbb P((U\cdot\nabla)v)(s)\,ds,
 \qquad U=v+c.
\]
Thus all entrance energy is regenerated by the nonlinearity.  Pairing
with \(v(\tau)\) and subtracting the instantaneous transport zero gives
an exact pressure-free heat-deviation work equal to
\(\norm{v(\tau)}_2^2\).

Writing \(M=\norm{U}_{L^{3,\infty}}\) and
\(Y=\norm{\nabla v}_2\), Lorentz heat smoothing forces a non-vanishing
half-order potential of \(MY\), failure of every finite temporal
\(L^{2,q}\) secondary index, and a quarter-order impulse of size
\(m_j\) at each q4 record.  An exact power--log ledger simultaneously
matches these requirements, the q4 \(M^2\) tail, the known \(Y^2\)
dissipation floor, and infinite amplitude-weighted dissipation.
Therefore the mild formulation and scalar norms do not close q4.  The
remaining task is a signed, orthogonal, or non-reusable decomposition of
the same-trajectory heat-deviation work.  All statements are conditional
repository results pending external review.  No Clay alternative is
proved.
\end{abstract}

\section{Epistemic status and setting}

This report assumes the full preceding q4 chain and proves no assertion
for arbitrary Clay-admissible data.  Let
\[
 v,c\in
 L^\infty(0,\tau_0;L^2_\sigma)
 \cap L^2(0,\tau_0;\dot H^1_\sigma)
\]
be smooth for every positive time and satisfy
\begin{equation}
\label{eq:entrance}
 \partial_\tau v-\nu\Delta v+
 \Pp((U\cdot\nabla)v)=0,
 \qquad U:=v+c,
 \qquad \nabla\cdot U=\nabla\cdot v=0.
\end{equation}
The boundary data are
\begin{align}
\label{eq:weak-zero}
 v(\tau)&\weakto0
 \quad\hbox{in }L^2
 \quad(\tau\downarrow0),\\
\label{eq:strong-c}
 c(\tau)&\longrightarrow c_0
 \quad\hbox{strongly in }L^2,\\
\label{eq:energy}
 \norm{v(\tau)}_2^2+
 2\nu\int_0^\tau\norm{\nabla v(s)}_2^2\,ds&=d>0.
\end{align}
In the full-defect construction \(c_0=-u_*\).

Set
\begin{equation}
\label{eq:MY}
 M(s):=\norm{U(s)}_{L^{3,\infty}},
 \qquad
 Y(s):=\norm{\nabla v(s)}_2,
 \qquad
 f(s):=M(s)Y(s).
\end{equation}
The q4 terminal tail and \eqref{eq:energy} give
\begin{equation}
\label{eq:L2-factors}
 M,Y\in L^2(0,\tau_0),
 \qquad
 f\in L^1(0,\tau_0).
\end{equation}
At exact records \(\tau_j\downarrow0\), the amplitude-inheritance theorem
gives
\begin{equation}
\label{eq:record-amplitude}
 \norm{v(\tau_j)}_{L^{3,\infty}}\gtrsim m_j.
\end{equation}

Write
\begin{equation}
\label{eq:heat}
 S(r):=e^{\nu r\Delta}.
\end{equation}
The heat semigroup commutes with \(\Pp\), preserves solenoidal fields,
and is self-adjoint on \(L^2\).

\section{The exact zero-boundary mild formula}

\begin{theorem}[Vanishing linear boundary datum]
\label{thm:mild}
For every fixed \(\tau\in(0,\tau_0)\),
\begin{equation}
\label{eq:mild}
 \boxed{
 v(\tau)
 =
 -\int_0^\tau
 S(\tau-s)\Pp((U\cdot\nabla)v)(s)\,ds
 }
\end{equation}
in \(L^2_\sigma(\R^3)\).
\end{theorem}

\begin{proof}
For \(0<\varepsilon<\tau\), ordinary positive-time mild evolution gives
\begin{equation}
\label{eq:positive-mild}
 v(\tau)
 =
 S(\tau-\varepsilon)v(\varepsilon)
 -
 \int_\varepsilon^\tau
 S(\tau-s)\Pp((U\cdot\nabla)v)(s)\,ds.
\end{equation}
Lorentz Holder yields
\begin{equation}
\label{eq:source}
 \norm{(U\cdot\nabla)v}_{L^{6/5,2}}
 \lesssim
 \norm{U}_{L^{3,\infty}}\norm{\nabla v}_2
 =f.
\end{equation}
The heat estimate
\begin{equation}
\label{eq:heat-L2}
 \norm{S(r)g}_2
 \lesssim
 (\nu r)^{-1/2}\norm{g}_{L^{6/5,2}}
\end{equation}
and \eqref{eq:L2-factors} show that the nonlinear integral in
\eqref{eq:positive-mild} converges strongly in \(L^2\) as
\(\varepsilon\downarrow0\).  Near \(s=0\), the heat time is bounded
away from zero and \(f\) is integrable.  Near \(s=\tau\), the fields are
smooth and the power \(r^{-1/2}\) is integrable.

For every \(g\in L^2_\sigma\),
\[
 \ip{S(\tau-\varepsilon)v(\varepsilon)}{g}
 =
 \ip{v(\varepsilon)}{S(\tau-\varepsilon)g}
 \longrightarrow0.
\]
Indeed, \eqref{eq:weak-zero} gives weak convergence of the first factor,
while \(S(\tau-\varepsilon)g\to S(\tau)g\) strongly.  Thus the linear
term in \eqref{eq:positive-mild} tends weakly to zero.  Since the other
term has a strong limit, the equality forces that weak limit to be the
zero strong limit.  Passing to the boundary proves \eqref{eq:mild}.
\end{proof}

\begin{remark}
This is not merely a distributional statement.  At every fixed positive
time, the heat-propagated escaping boundary oscillations vanish in
\(L^2\), and the complete field is supplied by the nonlinear integral.
\end{remark}

\section{Pressure-free heat-deviation work}

\begin{theorem}[Exact boundary work identity]
\label{thm:work}
For every \(\tau>0\),
\begin{equation}
\label{eq:work-dual}
 \boxed{
 \norm{v(\tau)}_2^2
 =
 -\int_0^\tau
 \ip{(U\cdot\nabla)v(s)}
 {S(\tau-s)v(\tau)-v(s)}\,ds
 }
\end{equation}
and equivalently
\begin{equation}
\label{eq:work-tensor}
 \boxed{
 \begin{aligned}
 \norm{v(\tau)}_2^2
 =
 \int_0^\tau\int_{\R^3}
 U_j(s)v_i(s)
 \partial_j\bigl(
 S(\tau-s)v_i(\tau)-v_i(s)
 \bigr)\,dx\,ds.
 \end{aligned}
 }
\end{equation}
The scalar time integrand is integrable when interpreted through the
first term in \eqref{eq:work-dual} and the pointwise cancellation below.
\end{theorem}

\begin{proof}
Pair \eqref{eq:mild} with \(v(\tau)\).  Self-adjointness of \(S\) and
the solenoidal range of \(S(\tau-s)v(\tau)\) remove the projection:
\begin{equation}
\label{eq:raw-work}
 \norm{v(\tau)}_2^2
 =
 -\int_0^\tau
 \ip{(U\cdot\nabla)v(s)}
 {S(\tau-s)v(\tau)}\,ds.
\end{equation}
For each \(s>0\),
\begin{equation}
\label{eq:transport-zero}
 \ip{(U\cdot\nabla)v(s)}{v(s)}=0
\end{equation}
because \(U\) is divergence free.  Subtract
\eqref{eq:transport-zero} inside \eqref{eq:raw-work}; this proves
\eqref{eq:work-dual}.  Spatial integration by parts proves
\eqref{eq:work-tensor}.

The raw integrand is absolutely integrable by
\eqref{eq:source}--\eqref{eq:heat-L2}.  Formula
\eqref{eq:work-dual} means that the pointwise scalar zero
\eqref{eq:transport-zero} is subtracted before time integration.  No
separate absolute estimate of the two spatial-gradient terms in
\eqref{eq:work-tensor} is claimed.
\end{proof}

\begin{remark}[Why the subtraction matters]
At \(s=\tau\), the heat deviation
\[
 S(\tau-s)v(\tau)-v(s)
\]
vanishes exactly.  A triangle inequality that estimates the two terms
separately loses this endpoint transport cancellation and introduces a
spurious singular upper bound.
\end{remark}

\begin{corollary}[Self and strong-drift split]
\label{cor:split}
The work decomposes as
\begin{equation}
\label{eq:split}
 \norm{v(\tau)}_2^2
 =
 \mathcal W_v(\tau)+\mathcal W_c(\tau),
\end{equation}
where, for \(z\in\{v,c\}\),
\begin{equation}
\label{eq:Wz}
 \mathcal W_z(\tau)
 :=
 \int_0^\tau\int_{\R^3}
 z_j(s)v_i(s)
 \partial_j\bigl(
 S(\tau-s)v_i(\tau)-v_i(s)
 \bigr)\,dx\,ds.
\end{equation}
Neither term has a known sign.
\end{corollary}

\begin{proof}
Both \(v\) and \(c\) are divergence free, so
\eqref{eq:transport-zero} holds separately with either field as the
advector.  Substitute \(U=v+c\) into \eqref{eq:work-tensor}.
\end{proof}

\begin{remark}
Strong convergence \(c(\tau)\to c_0\) controls the unamplified
\(L^2\) tail of \(c\), but gives no rate for the gradient of the
heat deviation in \eqref{eq:Wz}.  It therefore does not yet make
\(\mathcal W_c\) negligible.
\end{remark}

\section{Fractional impulse bounds}

For \(0<\alpha<1\), define the one-sided fractional potential
\begin{equation}
\label{eq:Ialpha}
 (I_\alpha f)(\tau)
 :=
 \int_0^\tau(\tau-s)^{\alpha-1}f(s)\,ds.
\end{equation}

\begin{theorem}[Half- and quarter-order entrance impulses]
\label{thm:impulse}
For every positive \(\tau\),
\begin{align}
\label{eq:half-upper}
 \norm{v(\tau)}_2
 &\lesssim
 \nu^{-1/2}I_{1/2}f(\tau),\\
\label{eq:quarter-upper}
 \norm{v(\tau)}_{L^{3,\infty}}
 &\lesssim
 \nu^{-3/4}I_{1/4}f(\tau).
\end{align}
Consequently, for every sufficiently small \(h>0\),
\begin{equation}
\label{eq:L1-floor}
 \boxed{
 \int_0^hM(s)Y(s)\,ds
 \gtrsim
 \sqrt{\nu d}\,h^{1/2}
 }
\end{equation}
and
\begin{equation}
\label{eq:product-floor}
 \boxed{
 \left(\int_0^hM(s)^2\,ds\right)
 \left(\int_0^hY(s)^2\,ds\right)
 \gtrsim
 \nu d\,h.
 }
\end{equation}
At every q4 record,
\begin{equation}
\label{eq:record-impulse}
 \boxed{
 I_{1/4}f(\tau_j)
 \gtrsim
 \nu^{3/4}m_j.
 }
\end{equation}
\end{theorem}

\begin{proof}
Equation \eqref{eq:half-upper} follows from
\eqref{eq:mild}, \eqref{eq:source}, and
\eqref{eq:heat-L2}.  The second heat estimate
\begin{equation}
\label{eq:heat-L3}
 \norm{S(r)g}_{L^{3,\infty}}
 \lesssim
 (\nu r)^{-3/4}\norm{g}_{L^{6/5,2}}
\end{equation}
proves \eqref{eq:quarter-upper}.

By \eqref{eq:energy},
\[
 \norm{v(\tau)}_2\longrightarrow\sqrt d.
\]
Thus \eqref{eq:half-upper} implies, for every sufficiently small
\(\tau\),
\begin{equation}
\label{eq:half-lower}
 I_{1/2}f(\tau)\gtrsim\sqrt{\nu d}.
\end{equation}
Integrate from zero to \(h\).  Tonelli's theorem gives
\[
\begin{aligned}
 \int_0^h I_{1/2}f(\tau)\,d\tau
 &=
 \int_0^h f(s)
 \int_s^h(\tau-s)^{-1/2}\,d\tau\,ds\\
 &=
 2\int_0^h f(s)(h-s)^{1/2}\,ds\\
 &\leq
 2h^{1/2}\int_0^h f(s)\,ds.
\end{aligned}
\]
Combining this with \eqref{eq:half-lower} proves
\eqref{eq:L1-floor}.  Cauchy--Schwarz proves
\eqref{eq:product-floor}.  Finally,
\eqref{eq:record-amplitude} and \eqref{eq:quarter-upper} prove
\eqref{eq:record-impulse}.
\end{proof}

\begin{corollary}[Failure of every finite secondary index]
\label{cor:lorentz}
For every \(1\leq q<\infty\),
\begin{equation}
\label{eq:no-lorentz}
 MY\notin L^{2,q}(0,\tau_0).
\end{equation}
\end{corollary}

\begin{proof}
Suppose \(f=MY\in L^{2,q}(0,\tau_0)\) for some finite \(q\).
The Lorentz norm is absolutely continuous on shrinking measurable sets,
and Lorentz Holder gives
\[
 \int_0^hf(s)\,ds
 \lesssim
 h^{1/2}\norm{f}_{L^{2,q}(0,h)}
 =
 o(h^{1/2}).
\]
This contradicts \eqref{eq:L1-floor}.
\end{proof}

\begin{remark}
The endpoint \(L^{2,\infty}\) is not excluded.  In fact
\eqref{eq:L1-floor} forces a non-vanishing critical tail rather than a
subcritical improvement.
\end{remark}

\begin{corollary}[q4 tail-product check]
\label{cor:q4-tail}
Suppose a canonical record obeys
\[
 \tau_j\asymp\frac{m_j^{-11/2}}j,
 \qquad
 \int_0^{\tau_j}M^2
 \lesssim\frac{m_j^{-7/2}}j.
\]
Then
\begin{equation}
\label{eq:weak-D}
 \int_0^{\tau_j}Y^2
 \gtrsim\nu d\,m_j^{-2}.
\end{equation}
\end{corollary}

\begin{proof}
Insert the two q4 scales into \eqref{eq:product-floor}.
\end{proof}

\begin{remark}
The preceding full-defect frequency theorem already gives the stronger
\(m_j^{-1}\)-scale dissipation tail.  Thus
\eqref{eq:weak-D} verifies consistency but does not close the branch.
\end{remark}

\section{Exact power--log saturation}

\begin{proposition}[Scalar endpoint ledger]
\label{prop:ledger}
For \(0<s<e^{-1}\), put
\[
 \ell(s):=\log(e/s),
\]
\begin{equation}
\label{eq:sharp-MY}
 M_\sharp(s):=
 s^{-2/11}\ell(s)^{-2/11},
 \qquad
 Y_\sharp(s):=
 s^{-9/22}\ell(s)^{1/11}.
\end{equation}
Then \(M_\sharp,Y_\sharp\in L^2(0,e^{-1})\), and as
\(h\downarrow0\),
\begin{align}
\label{eq:M-tail}
 \int_0^hM_\sharp^2\,ds
 &\asymp h^{7/11}\ell(h)^{-4/11},\\
\label{eq:Y-tail}
 \int_0^hY_\sharp^2\,ds
 &\asymp h^{2/11}\ell(h)^{2/11}.
\end{align}
Nevertheless,
\begin{equation}
\label{eq:weighted-divergence}
 \int_0^hM_\sharp Y_\sharp^2\,ds=\infty
 \qquad(h>0).
\end{equation}
If \(f_\sharp:=M_\sharp Y_\sharp\), then
\begin{align}
\label{eq:Ihalf-sharp}
 I_{1/2}f_\sharp(h)
 &\asymp
 h^{-1/11}\ell(h)^{-1/11},\\
\label{eq:Iquarter-sharp}
 I_{1/4}f_\sharp(h)
 &\asymp
 h^{-15/44}\ell(h)^{-1/11}.
\end{align}
\end{proposition}

\begin{proof}
The squares are
\[
 M_\sharp^2=s^{-4/11}\ell^{-4/11},
 \qquad
 Y_\sharp^2=s^{-9/11}\ell^{2/11},
\]
whose power exponents are strictly below one.  One-sided power--log
integration gives \eqref{eq:M-tail}--\eqref{eq:Y-tail}.
The weighted density is exactly
\[
 M_\sharp Y_\sharp^2=s^{-1},
\]
which proves \eqref{eq:weighted-divergence}.

Also
\[
 f_\sharp=s^{-13/22}\ell^{-1/11}.
\]
Scale \(s=h\rho\) in \eqref{eq:Ialpha}.  The beta integrals are finite
because \(13/22<1\), while regular variation permits replacement of
\(\ell(h\rho)\) by \(\ell(h)\) at leading order.  The powers are
\[
 1-\frac12-\frac{13}{22}=-\frac1{11},
 \qquad
 1-\frac34-\frac{13}{22}=-\frac{15}{44},
\]
which proves \eqref{eq:Ihalf-sharp}--\eqref{eq:Iquarter-sharp}.
\end{proof}

\begin{corollary}[Exact q4 record matching]
\label{cor:records}
Define
\begin{equation}
\label{eq:m-cont}
 m(h):=h^{-2/11}\ell(h)^{-2/11}.
\end{equation}
Then
\begin{equation}
\label{eq:quarter-room}
 \frac{I_{1/4}f_\sharp(h)}{m(h)}
 \asymp
 h^{-7/44}\ell(h)^{1/11}
 \longrightarrow\infty.
\end{equation}
If \(m_j\asymp2^{2j}\) and
\[
 h_j:=\frac{m_j^{-11/2}}j,
\]
then \(\ell(h_j)\asymp j\) and
\begin{align}
\label{eq:record-M}
 M_\sharp(h_j)&\asymp m_j,\\
\label{eq:record-Mtail}
 \int_0^{h_j}M_\sharp^2\,ds
 &\asymp\frac{m_j^{-7/2}}j,\\
\label{eq:record-Ytail}
 \int_0^{h_j}Y_\sharp^2\,ds
 &\asymp m_j^{-1}.
\end{align}
\end{corollary}

\begin{proof}
Equation \eqref{eq:quarter-room} is the quotient of
\eqref{eq:Iquarter-sharp} and \eqref{eq:m-cont}.  The remaining claims
follow by substituting \(h_j=m_j^{-11/2}/j\) and
\(\ell(h_j)\asymp j\) into
\eqref{eq:sharp-MY}, \eqref{eq:M-tail}, and
\eqref{eq:Y-tail}.
\end{proof}

\begin{remark}[Scope of the ledger]
The ledger simultaneously matches the q4 amplitude and \(M^2\) tail,
the known \(Y^2\) dissipation floor, infinite \(M\)-weighted
dissipation, the uniform half-order impulse, the record quarter-order
impulse, and failure of every finite secondary index.  It does not
construct spatial fields, a pressure, a Leray projection, signed work,
or one Navier--Stokes trajectory.
\end{remark}

\section{What changed in this round}

\subsection*{Robust conditional findings, subject to external review}

\begin{enumerate}
 \item The weak-zero boundary term vanishes in the exact \(L^2\) mild
 formula; all entrance energy is nonlinear.
 \item The complete positive energy is the pressure-free
 heat-deviation work \eqref{eq:work-dual}.
 \item The work retains the exact transport cancellation at its upper
 time endpoint.
 \item Every entrance forces the half-order tail
 \eqref{eq:L1-floor}, the product floor
 \eqref{eq:product-floor}, and failure of all finite temporal
 \(L^{2,q}\) secondary indices.
 \item The inherited q4 amplitude forces the recordwise quarter-order
 impulse \eqref{eq:record-impulse}.
 \item The power--log ledger satisfies every scalar consequence together
 with the prior q4 tail and dissipation laws.
\end{enumerate}

\subsection*{Closed shortcut}

Writing the entrance in mild form and using only the small terminal
\(L^2_\tau L^{3,\infty}_x\) norm of \(U\) cannot prove \(v=0\).  The
nonlinear term is necessarily order one at every small terminal time,
and scalar q4 data permit precisely the required endpoint
concentration.  Cauchy--Schwarz yields only
\eqref{eq:product-floor}, weaker than the existing spectral
dissipation floor.

\subsection*{Open obligations}

\begin{enumerate}
 \item Prove that the positive work
 \eqref{eq:work-dual} cannot be reused across an infinite nested family
 of q4 heat detectors.
 \item Derive a finite temporal \(L^{2,q}\) secondary index for
 \(MY\), with \(q<\infty\), from the same-trajectory coupling.
 \item Find a signed or orthogonal frequency decomposition of the work
 that forces a non-summable inverse-cascade charge.
 \item Use \(c(\tau)\to c_0\) strongly to make
 \(\mathcal W_c\) negligible at the moving clock without assuming an
 unavailable convergence rate.
 \item Combine the Hausdorff-one-null initial energy support with the
 heat-deviation identity to force a spatial-capacity charge.
 \item Subject the inherited chain and the present estimates to
 independent external mathematical review.
 \item Treat slower clocks, divergent normalised energy, R3B, and the
 other Clay alternatives separately.
\end{enumerate}

\begin{conjecture}[No reusable q4 boundary heat work]
\label{conj:nonreuse}
Under the exact same-trajectory energy-efficient q4 hypotheses, the
family
\[
 \mathcal W(\tau)
 :=
 -\int_0^\tau
 \ip{(U\cdot\nabla)v(s)}
 {S(\tau-s)v(\tau)-v(s)}\,ds
\]
cannot obey
\[
 \mathcal W(\tau)=\norm{v(\tau)}_2^2\longrightarrow d>0
\]
while the unweighted dissipation
\(\int_0^{\tau_0}\norm{\nabla v}^2_2\) remains finite.
\end{conjecture}

\begin{remark}
Conjecture~\ref{conj:nonreuse} is not proved.  The scalar ledger does not
test its sign, phase, projection, spatial capacity, or same-trajectory
non-reuse content.
\end{remark}

\section*{Conclusion}

The boundary formulation has become exact: no linear heat datum survives,
and the full defect is paid by a nonlinear heat-deviation work whose
endpoint transport cancellation is visible.  The associated scalar
impulse is necessarily critical and defeats every finite temporal
secondary index.  The q4 power--log ledger shows that these facts alone
remain compatible with finite unweighted dissipation.  Progress now
requires a genuinely new non-reuse or signed-frequency theorem.  The
Clay problem remains unsolved.

\end{document}
